\documentclass[11pt,a4paper]{article}
\usepackage[T1]{fontenc}
\usepackage[utf8]{inputenc}
\usepackage{lmodern,microtype}
\usepackage[a4paper,margin=25mm,headheight=15pt]{geometry}
\usepackage{amsmath,amssymb,amsthm,mathtools}
\usepackage{booktabs,longtable,array,enumitem}
\usepackage[dvipsnames]{xcolor}
\usepackage{fancyhdr,titlesec}
\usepackage[unicode,colorlinks=true,linkcolor=MidnightBlue,citecolor=MidnightBlue,urlcolor=MidnightBlue,breaklinks=true]{hyperref}
\usepackage{bookmark}
\usepackage{listings}
\definecolor{ink}{RGB}{27,49,70}
\titleformat{\section}{\large\bfseries\color{ink}}{\thesection}{0.7em}{}
\titleformat{\subsection}{\normalsize\bfseries\color{ink}}{\thesubsection}{0.7em}{}
\pagestyle{fancy}
\fancyhf{}
\fancyhead[L]{\small Countable exponential-feedback transseries}
\fancyhead[R]{\small Research draft}
\fancyfoot[C]{\thepage}
\renewcommand{\headrulewidth}{0.3pt}
\setlength{\parindent}{1.2em}
\setlength{\parskip}{3pt}
\setlength{\emergencystretch}{2em}
\setlist{itemsep=3pt,topsep=5pt,leftmargin=1.6em}
\numberwithin{equation}{section}
\newtheorem{theorem}{Theorem}[section]
\newtheorem{proposition}[theorem]{Proposition}
\newtheorem{lemma}[theorem]{Lemma}
\newtheorem{corollary}[theorem]{Corollary}
\theoremstyle{definition}
\newtheorem{definition}[theorem]{Definition}
\newtheorem{example}[theorem]{Example}
\theoremstyle{remark}
\newtheorem{remark}[theorem]{Remark}
\newcommand{\R}{\mathbb R}
\newcommand{\C}{\mathbb C}
\newcommand{\N}{\mathbb N}
\newcommand{\val}{\operatorname{ord}_q}
\newcommand{\Gev}{\mathcal G}
\newcommand{\Borel}{\mathcal B}
\newcommand{\supp}{\operatorname{supp}}
\newcommand{\lam}{\boldsymbol\lambda}
\newcommand{\mm}{\boldsymbol m}
\newcommand{\dd}{\mathrm d}
\newcommand{\ee}{\mathrm e}
\newcommand{\PhiL}{\Phi_{\lam}}
\newcommand{\Uhat}{\widehat U}
\newcommand{\Qhat}{\widehat Q}
\DeclareMathOperator*{\argmax}{arg\,max}
% ed. (2026-09-29): editorial-note environment for ProveIt cross-references;
% it uses no counter, so no author numbering changes.
\newenvironment{ednote}{\par\smallskip\begin{quote}\small\raggedright\textbf{Editorial note (ProveIt, 2026-09-29).}\ }{\end{quote}\smallskip}
\hypersetup{pdftitle={A sharp regularity classification for countable exponential-feedback transseries},pdfauthor={Research draft prepared with ChatGPT},pdfsubject={Convergence, exact Gevrey type, condensation, Borel growth and sectorial realization}}
\begin{document}
\begin{center}
{\LARGE\bfseries\color{ink} A sharp regularity classification for\par countable exponential-feedback transseries\par}
\vspace{0.6em}
{\large Convergence, exact Gevrey type, coefficient condensation,\par and sectorial realization}
\vspace{0.8em}

Research note prepared for Vladimir Reshetnikov\\
29 September 2026
\end{center}

\begin{abstract}
We study the formally well-defined equation
\[
 \Uhat(q)=\sum_{j\ge1}q^j\exp\!\bigl(\lambda_j\Uhat(q)\bigr),
 \qquad \lambda_j\ge0,
\]
as a model of countably many exponentially delayed corrections in transseries
inversion.  Formal existence alone imposes no growth restriction on the slopes
$\lambda_j$.  We derive a sharp separation between formal, analytic, and
Gevrey regularity.  The solution has positive radius of convergence exactly
when $\lambda_j=O(j)$.  More precisely, writing
$\Uhat=\sum_{n\ge1}u_nq^n$, for every $s>0$ we prove the exact identity
\[
 \limsup_{n\to\infty}\left(\frac{u_n}{(n!)^s}\right)^{1/n}
 =\left(\frac{s+1}{s}\right)^{s+1}
   \limsup_{j\to\infty}
      \frac{\lambda_j}{\bigl(j\log j\bigr)^{s+1}}.
\]
For $\lambda_j\sim a j^p(\log j)^\beta$, $p>1$, we determine
$\log u_n$ through its term of order $n$, and prove concentration of the
Lagrange weights on configurations with one large action of size
$\frac{p}{p-1}\frac{n}{\log n}$.  The quadratic model has an entire ordinary
Borel transform with $\log\log\Borel\Uhat(t)\sim2\sqrt t$, hence is not
Borel--Laplace summable on the positive ray.  Nevertheless, every nonnegative
slope sequence, even one of superpolynomial growth, admits a holomorphic
solution with this Poincar\'e expansion on proper subsectors of the left
half-plane.  We give explicit action-truncation bounds, transfer the Gevrey
classification to formal compositional inverses, and provide exact-arithmetic
checks and a focused research program.
\end{abstract}

\noindent\textbf{Status and scope.}
The results are proved here as conventional mathematical statements, not as
Lean theorems.  The sharp classification and its quantitative consequences
are the proposed contribution of this research draft.  Lagrange inversion,
Gevrey closure under reversion, and contraction arguments are classical
methods, not claimed inventions.  A priority claim against the entire
literature is not established, and no named longstanding conjecture is
claimed solved.

\clearpage
\setcounter{tocdepth}{1}
\hypersetup{bookmarksdepth=2}
\tableofcontents
\clearpage

\section{Problem selection and connection with ProveIt}
\label{sec:context}

\subsection{The repository boundary being investigated}

The ProveIt repository's consolidated \emph{Transseries: the
polynomial--logarithmic calculus and its inversions} connects formal
asymptotic scales, composition and reversal, Lambert-type equations,
analytic realization, and remainder certificates.  Its README explicitly
separates the human mathematics from the scope of the existing Lean APIs
\cite{repo}.  In particular, formal support control and analytic realization
are distinct tasks.  This paper investigates that distinction for a concrete
countable implicit kernel, rather than adding another finite list of
Lambert coefficients.

The repository snapshot consulted is commit
\texttt{04e06e032dff1966513bfba316966d52db3646b4}.  The relevant directory is
% ed. (2026-09-29): the pre-split directory quoted here is dead at the time of filing;
% replaced by its current location (the snapshot path is kept in the editorial note).
\begin{quote}\small\ttfamily
Analysis/Transseries/docs/\\
series-and-transseries/Transseries\_And\_Inversion/.
\end{quote}
% ed. (2026-09-29): current path of the canonical volume; the pinned snapshot's path recorded.
\begin{ednote}
The directory is given at its current location. At the pinned snapshot it was
\nolinkurl{Analysis/FabiusFunction/docs/semi-formalized-research-frontiers/drafts/series-and-transseries/Transseries_And_Inversion/};
the transseries material has since moved to \nolinkurl{Analysis/Transseries/}
(see \nolinkurl{docs/incoming/README.md}, batch 45). Packages filed later
beside this one in \nolinkurl{Analysis/Transseries/docs/series-and-transseries/}
answer several questions of Section~\ref{sec:future}; see the editorial notes there.
\end{ednote}
We use its documented program as motivation.  None of the proofs below
requires accepting an unverified frontier theorem from that repository.
Nor do we assert that the particular classification problem formulated here
was already posed there verbatim.

The central question is the following:
\begin{quote}
When countably many formally admissible exponential corrections feed back
into an implicit inverse, what growth of their feedback slopes is compatible
with convergence, with a specified Gevrey class, and with an actual
sectorial solution?
\end{quote}
A single-grid model already exhibits all three phenomena.  Thus the
obstruction does not depend on exotic supports or a failure of Hahn
summability.

\subsection{The model and its transseries interpretation}

Fix any sequence of finite real numbers $\lambda_j\ge0$.  Put
\begin{equation}
 \PhiL(q,u)=\sum_{j\ge1}q^j e^{\lambda_j u},
 \qquad \Uhat=\PhiL(q,\Uhat).
 \label{eq:kernel}
\end{equation}
Initially this is a \emph{formal} identity.  No analytic convergence of
$\PhiL$ near $(0,0)$ is assumed.  On writing $q=e^{-x}$, the solution becomes
\begin{equation}
 \Uhat(e^{-x})=\sum_{n\ge1}u_n e^{-nx}.
 \label{eq:transseries}
\end{equation}
This is a well-based, one-generator exponential transseries.  Each primitive
term has action $j$ and feedback slope $\lambda_j$.  Formal transseries
allow arbitrary real coefficients on such a support; coefficient growth
is a separate analytic question.  This distinction is part of the general
formal framework described by van der Hoeven \cite{vdh}.

The model is a nonlinear implicit inversion problem: one solves
$u-\PhiL(q,u)=0$ for $u$, and the resulting tangent-to-identity series itself
has a formal compositional inverse.  Section~\ref{sec:inverse} transfers the
classification to that inverse.

\subsection{What is classical and what is obtained here}

The coefficient formula in Section~\ref{sec:formal} is an application of
classical Lagrange inversion \cite{gessel}.  Infinite-colour inversion and
convergence theorems exist in much more general settings, including the
work of Jansen, Kuna, and Tsagkarogiannis \cite{jktcolors,jktvirial}.
We do not claim a first infinite-variable inversion theorem.

The specific contribution is the exact relationship between slope growth
and coefficient growth for \eqref{eq:kernel}, including the constant in the
Gevrey-type identity.  Its proof is an elementary but quantitative
optimization of the positive Lagrange formula.  The same argument produces
logarithmic coefficient asymptotics and a concentration theorem.  A second,
independent construction realizes the formal solution in left-hand sectors
without any slope-growth hypothesis.

Existing Gevrey implicit-function theorems require analytic or sectorial
uniformity hypotheses on the defining equation.  For example, Nikolaev's
theorem assumes suitable uniform Gevrey asymptotics in a genuine
neighborhood of the dependent variable \cite{nikolaev}.  Our rapidly growing
slopes prevent such a neighborhood at $(0,0)$.  The counterexamples below
therefore do not contradict those theorems.  Likewise, an entire Borel
transform is not by itself a Borel--Laplace summation theorem: growth at
infinity must also be controlled \cite{sauzin}.

\section{Main results and conventions}
\label{sec:main}

Throughout, $\log$ denotes the natural logarithm, $0^0=1$ in coefficient
formulas, and $\N$ includes zero.  All limits involving $\log j$ start at
$j\ge2$.  A multi-index $\mm=(m_1,m_2,\ldots)$ always has finite support.
Write
\[
 |\mm|=\sum_jm_j,\qquad w(\mm)=\sum_jjm_j,\qquad
 \lambda\cdot\mm=\sum_j\lambda_jm_j,\qquad
 \mm!=\prod_jm_j!.
\]

\begin{definition}
For $s\ge0$, a formal series $\widehat f=\sum f_nq^n$ is in $\Gev_s$ if
there are $C,A>0$ such that $|f_n|\le CA^n(n!)^s$ for every $n$.
For $s>0$ its factorial-normalized type is
\[
 T_s(\widehat f)=\limsup_{n\to\infty}
       \left(\frac{|f_n|}{(n!)^s}\right)^{1/n}.
\]
Types may be zero or infinite.  The class $\Gev_0$ is precisely the class
of convergent power series.
\end{definition}

\begin{theorem}[Sharp regularity classification]
\label{thm:classification}
Equation~\eqref{eq:kernel} has a unique solution
$\Uhat\in q\R[[q]]$, with $u_1=1$ and $u_n\ge1$.
Its radius of convergence is positive if and only if
\begin{equation}
 \sup_{j\ge1}\frac{\lambda_j}{j}<\infty.
 \label{eq:analyticcriterion}
\end{equation}
For every $s>0$, the following exact identity holds in $[0,\infty]$:
\begin{equation}
 \boxed{\quad
 T_s(\Uhat)=C_s A_s(\lam),\qquad
 C_s=\left(\frac{s+1}{s}\right)^{s+1},\qquad
 A_s(\lam)=\limsup_{j\to\infty}
  \frac{\lambda_j}{(j\log j)^{s+1}}.
 \quad}
 \label{eq:exacttype}
\end{equation}
Consequently,
\begin{equation}
 \Uhat\in\Gev_s
 \quad\Longleftrightarrow\quad
 \lambda_j=O\!\left((j\log(j+1))^{s+1}\right),
 \qquad s>0.
 \label{eq:gevcriterion}
\end{equation}
\end{theorem}

There is a real endpoint distinction: \eqref{eq:gevcriterion} must not be
extrapolated to $s=0$.  For example, $\lambda_j=j\log(j+1)$ gives a divergent
series that belongs to every positive Gevrey class.

\begin{theorem}[Regularly varying slopes]
\label{thm:regular}
Suppose, for constants $a>0$, $p>1$, and $\beta\in\R$, that
\[
 \lambda_j\sim a j^p(\log j)^\beta.
\]
Then, as $n\to\infty$,
\begin{align}
 \log u_n={}&(p-1)n\log n+(\beta-p)n\log\log n
 \notag\\
 &+\left[\log a+p\log\frac{p}{p-1}-(p-1)\right]n+o(n).
 \label{eq:logasymptotic}
\end{align}
Under the probability distribution on $w(\mm)=n$ proportional to
\[
 W(\mm)=\frac{(\lambda\cdot\mm)^{|\mm|-1}}{\mm!},
\]
let $K=|\mm|$, $R_n=n-K+1$, and
$J_{\max}=\max\{j:m_j>0\}$.  Then
\begin{equation}
 \frac{R_n\log n}{n}\longrightarrow\frac{p}{p-1},
 \qquad
 \frac{J_{\max}}{R_n}\longrightarrow1
 \label{eq:condensationintro}
\end{equation}
in probability.  Each fixed positive deviation in either statement has
probability bounded by $\exp(-c n+o(n))$ for some $c>0$.
\end{theorem}

The concentration statement concerns precisely this finite Lagrange-weight
probability space.  It is not a blanket assertion about all models of
random rooted trees.

\begin{theorem}[An entire Borel transform with excessive growth]
\label{thm:borelintro}
For $\lambda_j\sim a j^p$, $p>1$, set $s=p-1$ and define
\[
 B_s(t)=\sum_{n\ge1}\frac{u_n}{(n!)^s}t^n.
\]
Then $B_s$ is entire and
\begin{equation}
 \lim_{t\to+\infty}\frac{\log\log B_s(t)}{t^{1/p}}
 =a^{1/p}\frac{p}{p-1}.
 \label{eq:doublelog}
\end{equation}
In particular, for $\lambda_j=j^2$ the ordinary integrated Borel transform
$B_1$ satisfies $\log\log B_1(t)\sim2\sqrt t$.  The ordinary positive-ray
Laplace integral of $B_1$ diverges for every positive value of $q$.
\end{theorem}

\begin{theorem}[Universal left-sector realization]
\label{thm:sectorintro}
For every sequence $\lambda_j\ge0$ and every $\eta\in(0,1)$, there is
$r_0>0$ such that the sector
\[
 S_{\eta,r_0}=\{q:0<|q|<r_0,\ \Re q<-\eta|q|\}
\]
supports a holomorphic solution $U$ of the literal, absolutely convergent
identity $U=\PhiL(q,U)$, with $U=q+O(q^2)$.  It has Poincar\'e expansion
$\Uhat$ on that sector.  Its finite-action approximations satisfy an
explicit error bound of order $|q|^{N+1}$, uniformly in the number $N$ of
retained actions.
\end{theorem}

Theorems~\ref{thm:classification}--\ref{thm:sectorintro} address different
questions.  Sectorial existence does not assert convergence at the origin;
Poincar\'e asymptotics does not assert uniform Gevrey remainder bounds;
and an entire Borel transform does not assert adequate Laplace growth.

\section{Formal existence and the positive coefficient formula}
\label{sec:formal}

\begin{lemma}[Admissibility and coefficientwise contraction]
\label{lem:formal}
For any $u\in q\R[[q]]$, the right-hand side of \eqref{eq:kernel} is a
well-defined element of $q\R[[q]]$.  If $u-v$ has order $m\ge1$, then
\[
 \val\bigl(\PhiL(q,u)-\PhiL(q,v)\bigr)\ge m+1.
\]
Consequently the iterates $U^{[0]}=0$ and
$U^{[h+1]}=\PhiL(q,U^{[h]})$ stabilize coefficientwise, with the coefficient
of $q^n$ stable once $h\ge n$.
\end{lemma}

\begin{proof}
At coefficient $q^n$, only $j\le n$ can occur.  In the exponential attached
to such a $j$, only powers $u^k$ with $k\le n-j$ contribute.  Thus every
coefficient is a finite sum.  The formal difference
$e^{\lambda_j u}-e^{\lambda_j v}$ is divisible by $u-v$; multiplication by
$q^j$ raises order by at least one.  This proves the contraction estimate.
Completeness of $\R[[q]]$ in the coefficientwise, or $q$-adic, topology gives
a fixed point.  If two fixed points differed first in degree $m$, the same
estimate would make their difference have order at least $m+1$, a
contradiction.  Iteration also proves the stabilization assertion.
\end{proof}

\begin{proposition}[Two equivalent Lagrange formulas]
\label{prop:coefficients}
The coefficients of $\Uhat$ are
\begin{align}
 u_n
 &=\sum_{\substack{\mm\ge0\\w(\mm)=n}}
       \frac{(\lambda\cdot\mm)^{|\mm|-1}}{\mm!},
 \label{eq:partitionformula}\\
 &=\sum_{k=1}^n\frac1{k!}
       \sum_{\substack{j_1+\cdots+j_k=n\\j_i\ge1}}
       (\lambda_{j_1}+\cdots+\lambda_{j_k})^{k-1}.
 \label{eq:compositionformula}
\end{align}
In particular, $u_n\ge1$, and each $u_n$ is nondecreasing in every slope.
\end{proposition}

\begin{proof}
Introduce a marker $t$ for the number of primitive factors and solve
$V=t\PhiL(q,V)$.  Classical formal Lagrange inversion \cite{gessel} gives
\[
 [t^k]V=\frac1k[z^{k-1}]\PhiL(q,z)^k.
\]
For the ordered tuple $(j_1,\ldots,j_k)$, the product of exponentials is
$\exp((\sum_i\lambda_{j_i})z)$, whose coefficient of $z^{k-1}$ is
$(\sum_i\lambda_{j_i})^{k-1}/(k-1)!$.  This proves
\eqref{eq:compositionformula}.  Setting $t=1$ is legitimate coefficientwise
because $k\le n$ in degree $q^n$.  Grouping ordered tuples by their
multiplicities gives $k!/\mm!$ tuples in each group and proves
\eqref{eq:partitionformula}.  The tuple $(n)$ contributes one, and all
terms are nonnegative.  Monotonicity is immediate.
\end{proof}

\subsection{A useful lower bound}
For integers $r\ge2$ and $k\ge1$ with $r+k-1=n$, retain the multiplicities
$m_r=1$, $m_1=k-1$.  Equation~\eqref{eq:partitionformula} yields
\begin{equation}
 u_n\ge
 \frac{(\lambda_r+(k-1)\lambda_1)^{k-1}}{(k-1)!}
 \ge\frac{\lambda_r^{k-1}}{(k-1)!}.
 \label{eq:singlelarge}
\end{equation}
This one-large-action contribution will determine the sharp exponential
scale.  The case $r=1$ must be treated separately: then there is only one
multiplicity vector, not $k$ distinct placements of a distinguished action.

\subsection{Finite truncations}
Let $\Uhat_N$ solve
\begin{equation}
 \Uhat_N=\sum_{j=1}^Nq^j e^{\lambda_j\Uhat_N}.
 \label{eq:finitekernel}
\end{equation}
Its defining function is analytic near $(0,0)$, and its derivative with
respect to the left-hand unknown at the origin is one.  Thus $\Uhat_N$
is a convergent germ.  The formal locality above gives
\begin{equation}
 [q^n]\Uhat_N=[q^n]\Uhat\qquad(n\le N).
 \label{eq:finiteagreement}
\end{equation}
It is precisely the lack of a uniform analytic neighborhood as $N$ grows
that allows the infinite formal limit to diverge.

\section{The exact convergence criterion and positive certificates}
\label{sec:convergence}

\begin{proposition}[Convergence at a positive point]
\label{prop:majorant}
For $0<q<1$, the nonnegative series $\sum_nu_nq^n$ is finite if and only if
there exists $r>0$ such that
\begin{equation}
 \sum_{j\ge1}q^j e^{\lambda_jr}\le r.
 \label{eq:supersolution}
\end{equation}
When finite, its sum is the least nonnegative fixed point of
$r\mapsto\PhiL(q,r)$.
\end{proposition}

\begin{proof}
Assume \eqref{eq:supersolution}.  Monotone real iteration from zero remains
in $[0,r]$.  Its coefficients are those of the formal iterates in
Lemma~\ref{lem:formal}, are nonnegative, and increase to $u_n$.
The monotone convergence theorem, applied to these nonnegative series,
gives $\sum_nu_nq^n\le r$.  Applying the same theorem to the exponential
expansions verifies the fixed-point identity.  Iteration from zero is
below every nonnegative fixed point, proving minimality.

Conversely, if $u=\sum_nu_nq^n<\infty$, substitute the sum into the formal
identity and use Tonelli's theorem for the nonnegative terms.  It gives
$u=\sum_jq^je^{\lambda_ju}$.  Since $u\ge q>0$, the value $r=u$ is an
admissible supersolution.
\end{proof}

\begin{proof}[Proof of the convergence part of Theorem~\ref{thm:classification}]
If $\lambda_j\le Mj$ for all $j$, then for sufficiently small $|q|$ and
bounded $|u|$,
\[
 \sum_j |q|^j e^{\lambda_j|u|}
 \le\sum_j\bigl(|q|e^{M|u|}\bigr)^j.
\]
Hence $\PhiL$ is holomorphic in a product neighborhood of $(0,0)$.
The analytic implicit-function theorem applies to $u-\PhiL(q,u)$ because
its $u$-derivative at the origin is one.

If $\sup_j\lambda_j/j=\infty$, take any $0<q<1$ and $r>0$.
Along a subsequence,
\[
 j\log q+\lambda_jr
 =j\left(\log q+r\frac{\lambda_j}{j}\right)\longrightarrow+\infty.
\]
The summands in \eqref{eq:supersolution} do not even tend to zero.
Proposition~\ref{prop:majorant} excludes convergence at every positive
$q$.  Nonnegativity of the coefficients then implies radius zero.
\end{proof}

\begin{remark}
The divergence proof is not merely an upper-majorant failure.  Positivity
makes the supersolution criterion necessary as well as sufficient.
For signed or complex coefficients this equivalence generally disappears.
\end{remark}

\subsection{Quantitative certificates inside a convergence domain}
Suppose, for $Q,r>0$, that
\begin{equation}
 \PhiL(Q,r)\le r,
 \qquad
 \kappa=\partial_u\PhiL(Q,r)
       =\sum_jQ^j\lambda_j e^{\lambda_jr}<1.
 \label{eq:contractcertificate}
\end{equation}
Then the fixed-point map is a contraction on $|u|\le r$ for $|q|\le Q$.
The infinite solution and the solution $U_N$ of the first $N$ actions obey
\begin{equation}
 |U(q)-U_N(q)|\le
 \frac{\sum_{j>N}Q^j e^{\lambda_jr}}{1-\kappa},
 \qquad |q|\le Q.
 \label{eq:positivetail}
\end{equation}
This follows by subtracting the two fixed-point equations and bounding
the derivative on the line segment joining their values.

A different certificate bounds truncation by total action rather than by
primitive action.  If merely $\PhiL(Q,r)\le r$, then for $|q|<Q$,
\begin{equation}
 \left|U(q)-\sum_{n=1}^Nu_nq^n\right|
 \le r\left(\frac{|q|}{Q}\right)^{N+1}.
 \label{eq:taylortail}
\end{equation}
Indeed, factor $(|q|/Q)^{N+1}$ from the nonnegative coefficient tail and use
$\sum_nu_nQ^n\le r$.  Equations~\eqref{eq:positivetail} and
\eqref{eq:taylortail} certify different truncation operations; they should
not be interchanged without checking their hypotheses.

\section{A coefficient-envelope lemma}
\label{sec:envelope}

This section contains the quantitative step behind both the asymptotic
formula and the exact Gevrey type.  The details matter: replacing the number
of compositions by $2^n$ throughout would lose a constant of order $n$ and
would not prove \eqref{eq:exacttype}.

Fix $p>1$, $\beta\in\R$, and write
\[
 s=p-1,\qquad f(r)=r^p\bigl(\log(r+1)\bigr)^\beta.
\]
For $A>0$ define
\begin{equation}
 X_n(A)=s n\log n+(\beta-p)n\log\log n
       +\left[\log A+p\log\frac p s-s\right]n.
 \label{eq:Xn}
\end{equation}

\begin{lemma}[Sharp logarithmic envelope]
\label{lem:envelope}
For every $A>0$ and $D\ge0$,
\begin{align}
 \max_{2\le r\le n}
 \log\frac{(A f(r))^{n-r}}{(n-r)!}
 &=X_n(A)+o(n),
 \label{eq:lowerenvelope}\\
 \log\sum_{r=1}^n
 \binom{n-1}{r-1}
 \frac{\bigl(A f(r)+D(n-r+1)\bigr)^{n-r}}{(n-r+1)!}
 &=X_n(A)+o(n).
 \label{eq:upperenvelope}
\end{align}
The relevant values of $r$ satisfy $r\log n/n\to p/s$ at the
exponential scale.  More precisely, restricting the upper sum to
$|r\log n/n-p/s|\ge\varepsilon$ lowers its logarithm by
$c_\varepsilon n+o(n)$ for some $c_\varepsilon>0$.
\end{lemma}

\begin{proof}
Put $L=\log n$, $k=n-r+1$, and
$M=A f(r)+Dk$.  The logarithm of an upper-sum term is
\begin{equation}
 E_n(r)=\log\binom{n-1}{r-1}+(k-1)\log M-\log k!.
 \label{eq:En}
\end{equation}
All occurrences of $M$ have polynomial growth in $n$ times a power of
$\log n$, and $M$ is bounded away from zero.  Uniform Stirling estimates
therefore permit errors $O(\log n)$ when replacing factorials by their
leading logarithms.

\paragraph{A lower comparison.}
Choose
\[
 r=\left\lfloor\frac p s\frac nL\right\rfloor.
\]
Then $r\to\infty$, $r/n\to0$, and
\[
 \log f(r)=pL+(\beta-p)\log L+p\log(p/s)+o(1).
\]
Substituting this into the lower expression and using Stirling gives
$X_n(A)+o(n)$.  Since the corresponding upper term differs from this
lower term by the factor $\binom{n-1}{r-1}/k$ and by replacing
$Af(r)$ with a larger quantity, its logarithm is also at least
$X_n(A)+o(n)$.  Here $\log k=o(n)$.

\paragraph{Localization of any maximizing row.}
For $r\ge\delta n$ with fixed $\delta>0$, we have
$k\le(1-\delta)n+1$ and
\[
 E_n(r)\le(1-\delta)s nL+O(n\log L)+O(n).
\]
The entropy bound $\log\binom{n-1}{r-1}\le n\log2$ is sufficient
\emph{at this preliminary step}.  Such rows are smaller than the lower
comparison by a positive multiple of $nL$.  A maximizing row thus has
$r/n\to0$.  For these rows
\begin{equation}
 \log\binom{n-1}{r-1}=o(n).
 \label{eq:smallentropy}
\end{equation}
Write $\delta_n=r/n$.  Comparing with the lower value now forces
$\log M\ge pL-O(\log L)$.  Consequently $M/n\to\infty$, the term $Dk$
is negligible in $M$, and
\[
 \log r=L-O(\log L),\qquad
 \log\log(r+1)=\log L+o(1).
\]
At such a row, Stirling and \eqref{eq:smallentropy} give
\begin{align*}
 \frac{E_n(r)}n
  &=(1-\delta_n)\bigl[sL+p\log\delta_n+\beta\log L
             +\log A-\log(1-\delta_n)+1\bigr]+o(1).
\end{align*}
The same comparison, using $\log\delta_n\le0$, implies
$\delta_n L=O(\log L)$.  Hence
$\delta_n\log L=o(1)$ and $\delta_n\log\delta_n=o(1)$.
The previous display reduces to
\begin{equation}
 \frac{E_n(r)}n
 =sL+(\beta-p)\log L+\log A+1
      +p\log y-sy+o(1),
 \qquad y=\frac{rL}{n}.
 \label{eq:saddlefunctional}
\end{equation}

\paragraph{The one-dimensional maximum.}
The function
\[
 g(y)=p\log y-sy\qquad(y>0)
\]
is strictly concave, tends to $-\infty$ at both ends, and has its unique
maximum at $y_*=p/s$.  Its maximum value is $p\log(p/s)-p$.
Substitution into \eqref{eq:saddlefunctional} proves that the largest
upper term has logarithm at most $X_n(A)+o(n)$.  A sum of $n$ terms changes
the logarithm by at most $\log n=o(n)$, proving
\eqref{eq:upperenvelope}.  The upper expression also bounds every lower
expression up to the harmless factor $k$, because the binomial factor is
at least one.  This proves \eqref{eq:lowerenvelope}.

For the restricted sum, a contradiction argument gives the same
localization for any sequence of rows whose values are within $O(n)$ of
the maximum.  On the closed set $|y-y_*|\ge\varepsilon$, strict concavity
and the behavior at the ends give
$g(y)\le g(y_*)-c_\varepsilon$.  Rows outside the localized range have
already lost more than a fixed multiple of $n$.  Applying
\eqref{eq:saddlefunctional} proves the final assertion.  This argument also
makes the $o(n)$ estimates uniform over the rows needed for that assertion.
\end{proof}

\subsection{From eventual convexity to a uniform upper bound}

\begin{lemma}[Convex majorant]
\label{lem:convex}
Suppose $\lambda_j\sim a f(j)$.  For every $A>a$, there is $D\ge0$ such
that every composition $j_1+\cdots+j_k=n$, $j_i\ge1$, satisfies
\begin{equation}
 \lambda_{j_1}+\cdots+\lambda_{j_k}
 \le A f(n-k+1)+Dk.
 \label{eq:convexmajorant}
\end{equation}
\end{lemma}

\begin{proof}
Direct differentiation gives
$f''(t)\sim p(p-1)t^{p-2}(\log t)^\beta>0$.
Thus $f$ is convex and increasing eventually.  Extend its tail to a convex
function $h$ on $[1,\infty)$ by continuing a tangent line to the left.
Then $h=f$ on a tail.  Since $A>a$, increasing a constant $D_0$ handles
the finitely many exceptional indices and gives
$\lambda_j\le Ah(j)+D_0$ for all $j$.

For a convex function, shifting mass from the smaller of two arguments
larger than one to the larger cannot decrease their sum.  Repeating gives
\[
 \sum_{i=1}^kh(j_i)\le h(n-k+1)+(k-1)h(1).
\]
On the finite interval where $h\ne f$, their difference is bounded.
Absorbing this bound and $Ah(1)+D_0$ into a nonnegative $D$ proves the
claim.  This also explains why arbitrarily large but finitely many initial
slopes do not alter the asymptotic constant.
\end{proof}

\section{Logarithmic asymptotics and concentration}
\label{sec:asymptotics}

\begin{proof}[Proof of the coefficient part of Theorem~\ref{thm:regular}]
For any $A>a$, apply Lemma~\ref{lem:convex} in
\eqref{eq:compositionformula}.  There are
$\binom{n-1}{k-1}=\binom{n-1}{r-1}$ compositions with $k=n-r+1$ parts.
Lemma~\ref{lem:envelope} therefore gives
\[
 \log u_n\le X_n(A)+o(n).
\]
For a lower bound, choose $r=\lfloor (p/s)n/\log n\rfloor$ in
\eqref{eq:singlelarge}.  Since $\lambda_r/(a f(r))\to1$, the lower
calculation in Lemma~\ref{lem:envelope} gives
\[
 \log u_n\ge X_n(a)+o(n).
\]
Letting $A\downarrow a$ proves \eqref{eq:logasymptotic}.
\end{proof}

\subsection{Concentration of the total excess action}

The integer $n-K=\sum_j(j-1)m_j$ is the total action remaining after giving
one unit to each of the $K$ factors.  It is a more useful variable than
$K$ itself: almost all factors have unit action on the leading scale.

\begin{proposition}
\label{prop:excess}
Under the assumptions of Theorem~\ref{thm:regular}, for every
$\varepsilon>0$ there is $c_\varepsilon>0$ such that
\[
 \Pr\!\left(\left|\frac{(n-K+1)\log n}{n}-\frac p{p-1}\right|
                \ge\varepsilon\right)
 \le\exp(-c_\varepsilon n+o(n)).
\]
\end{proposition}

\begin{proof}
The total unnormalized weight at a fixed value of $K=k$ is exactly the
$k$-row of \eqref{eq:compositionformula}.  Bound the rows by
Lemma~\ref{lem:convex}.  The restricted-envelope conclusion in
Lemma~\ref{lem:envelope} gives a strict exponential loss outside the stated
window.  Choose $A>a$ sufficiently close to $a$ that the increase
$n\log(A/a)$ is smaller than half that loss.  Divide by the lower estimate
$u_n=\exp(X_n(a)+o(n))$.
\end{proof}

\subsection{Concentration on one large action}

\begin{proposition}
\label{prop:giant}
Under the same assumptions, for every $\varepsilon>0$ there is
$c_\varepsilon>0$ such that
\[
 \Pr\!\left(J_{\max}\le(1-\varepsilon)(n-K+1)\right)
 \le\exp(-c_\varepsilon n+o(n)).
\]
\end{proposition}

\begin{proof}
It suffices to consider $0<\varepsilon<1/2$, since increasing
$\varepsilon$ makes the event smaller.  By Proposition~\ref{prop:excess},
we may restrict to
$r=n-k+1$ in a fixed compact multiple of $n/\log n$, losing only an
exponentially small probability.  In this range $k\sim n$ and
$f(r)/n\to\infty$.

Set $b=\lfloor(1-\varepsilon)r\rfloor$.  For a convex function $h$ as in
Lemma~\ref{lem:convex}, maximize $\sum_i h(j_i)$ subject to
$1\le j_i\le b$ and $\sum_i(j_i-1)=r-1$.  Since $b>r/2$ for large $r$,
a maximizing configuration has at most two entries greater than one:
one at $b$, the other at $r-b+1$.  This follows by the same pairwise
mass-transfer argument, stopping when an entry reaches the cap.  Thus
\[
 \sum_i\lambda_{j_i}
 \le A\{f(b)+f(r-b+1)\}+D'k.
\]
Uniformly in the restricted range of $r$,
\[
 \frac{f(b)+f(r-b+1)}{f(r)}
 \longrightarrow d_\varepsilon
 :=(1-\varepsilon)^p+\varepsilon^p<1.
\]
The $D'k$ term is negligible.  Choose $A>a$ close enough that
$(A/a)d_\varepsilon<1$.  Compared with the unrestricted row bound, the
capped slope sum has a fixed multiplicative loss.  Its power is $k-1\sim n$,
so the row loses $c n+o(n)$ in its logarithm.  The composition-count entropy
is $o(n)$ in this range and cannot repair the loss.  Summing the at most
$n$ rows and dividing by $u_n=\exp(X_n(a)+o(n))$ completes the proof.
\end{proof}

Together these propositions finish Theorem~\ref{thm:regular}.  They justify
the one-large-action mechanism at exponential precision.  They do
\emph{not} give a local limit theorem, a fluctuation scale about the true
finite-$n$ center, or a full multiplicative asymptotic for $u_n$.
% ed. (2026-09-29): cross-reference to later packages that close these gaps.
\begin{ednote}
For power slopes these gaps are closed by later packages in this series;
see the editorial note after the question ``Beyond logarithmic coefficient
asymptotics'' in Section~\ref{sec:future}.
\end{ednote}

\section{The exact Gevrey-type identity}
\label{sec:gevrey}

\begin{proof}[Proof of \eqref{eq:exacttype}]
Fix $s>0$ and put $p=s+1$.  Denote $A_s(\lam)$ by $A_*$.  First suppose
$A_*<\infty$.  For every $A>A_*$ there is $D\ge0$ such that
\[
 \lambda_j\le\widetilde\lambda_j
 :=A\bigl(j\log(j+1)\bigr)^p+D
\]
for every $j$.  Coefficientwise monotonicity gives
$u_n(\lam)\le u_n(\widetilde\lam)$.  Apply
Theorem~\ref{thm:regular} to $\widetilde\lambda_j$ with $\beta=p$:
\[
 \log u_n(\widetilde\lam)
 =s n\log n+\left[\log A+p\log(p/s)-s\right]n+o(n).
\]
Since $\log n!=n\log n-n+o(n)$,
\[
 T_s(\Uhat)\le A(p/s)^p.
\]
Let $A\downarrow A_*$.  This also proves the upper bound when $A_*=0$.

For the lower bound, take any $b>0$ with $b<A_*$; when $A_*=\infty$,
$b$ is arbitrary.  Infinitely many $j$ satisfy
\[
 \lambda_j\ge b(j\log j)^p.
\]
Along these indices, choose
\[
 k=\left\lfloor\frac s p j\log j\right\rfloor,
 \qquad n=j+k.
\]
The multiplicities $m_j=1$ and $m_1=k$ give
$u_n\ge\lambda_j^k/k!$.  To record the sharp constant explicitly, put
$c=s/p$ and $L=\log j$.  Then
\[
 k=cjL+O(1),\quad n=cjL+j+O(1),\quad
 \log n=L+\log L+\log c+o(1).
\]
Stirling's formula, with the relative difference $j/n$ retained, gives
\begin{align}
 \frac1n\log\frac{\lambda_j^k}{k!(n!)^s}
 &\ge \log b-p\log c+p-\frac s c+o(1)
 \notag\\
 &=\log b+p\log(p/s)+o(1).
 \label{eq:typelowercalc}
\end{align}
Therefore $T_s(\Uhat)\ge b(p/s)^p$.  Letting $b\uparrow A_*$ proves the
lower bound, including infinity.  When $A_*=0$, the lower bound is simply
nonnegativity.  This proves the identity in every case.
\end{proof}

\subsection{Consequences and endpoint behavior}

\begin{corollary}[Borel radii and zero type]
\label{cor:borelradius}
For $s>0$, the factorial-normalized transform
$\sum_nu_nt^n/(n!)^s$ has radius
\begin{equation}
 R_s=\frac{1}{C_s A_s(\lam)},
 \label{eq:borelradius}
\end{equation}
with $1/0=\infty$ and $1/\infty=0$.  In particular it is entire exactly
when $\lambda_j=o((j\log j)^{s+1})$.
For the normalization $\Gamma(1+sn)$, the radius is instead
$s^s/(C_sA_s(\lam))$.
\end{corollary}

\begin{proof}
The first statement is Cauchy--Hadamard and
Theorem~\ref{thm:classification}.  Stirling gives
$\Gamma(1+sn)^{1/n}/(n!)^{s/n}\to s^s$, which proves the last assertion.
\end{proof}

\begin{corollary}[The least Gevrey index]
\label{cor:index}
Let
\[
 \rho=\limsup_{j\to\infty}\frac{\log(1+\lambda_j)}{\log j}.
\]
Then
\begin{equation}
 \inf\{s\ge0:\Uhat\in\Gev_s\}=\max(0,\rho-1),
 \label{eq:index}
\end{equation}
with value infinity when $\rho=\infty$.  Membership at the infimum is
not automatic.
\end{corollary}

\begin{proof}
For $s>\max(0,\rho-1)$, a positive power of $j$ separates
$\lambda_j$ from $(j\log j)^{s+1}$ eventually, so
\eqref{eq:gevcriterion} applies.  If $0<s<\rho-1$, the reverse separation
holds on a subsequence and excludes $\Gev_s$.  The case $s=0$ is covered
by the convergence criterion.  These implications prove the infimum,
without making an unwarranted assertion about its attainment.
\end{proof}

For $\lambda_j\sim a j^p(\log j)^\beta$ with $p>1$, the critical index is
$s=p-1$.  At that index the type is zero when $\beta<p$, is
$a(p/(p-1))^p$ when $\beta=p$, and is infinite when $\beta>p$.
Thus the logarithmic factor decides whether the critical Gevrey class is
attained, and whether its Borel radius is infinite, finite, or zero.

\begin{example}[Sparse spikes are detected sharply]
Fix $s>0$ and $a>0$, and let
\[
 \lambda_j=\begin{cases}
 a(j\log j)^{s+1},&j=2^{2^m},\ m\ge1,\\
 0,&\text{otherwise}.
 \end{cases}
\]
Then $T_s(\Uhat)=aC_s$.  The density of the large slopes is irrelevant to
this limsup statement.  We do not infer a full coefficient asymptotic
from this sparse example; the regular-variation hypothesis of
Theorem~\ref{thm:regular} does not hold.
\end{example}

\section{Borel growth and the positive-direction obstruction}
\label{sec:borel}

\subsection{An elementary entire-function lemma}

\begin{lemma}
\label{lem:entiregrowth}
Let $b_n\ge0$ for all $n$, with $b_n>0$ eventually, and suppose, for $p>0$ and $D\in\R$, that
\begin{equation}
 \log b_n=-p n\log\log n+Dn+o(n).
 \label{eq:bnassumption}
\end{equation}
Then $B(t)=\sum_{n\ge1}b_nt^n$ is entire and
\begin{equation}
 \lim_{t\to+\infty}\frac{\log\log B(t)}{t^{1/p}}=e^{D/p}.
 \label{eq:entiregrowth}
\end{equation}
The same conclusion holds after changing finitely many coefficients to
arbitrary nonnegative values.
\end{lemma}

\begin{proof}
Equation~\eqref{eq:bnassumption} gives $b_n^{1/n}\to0$, so $B$ is entire.
Put $c=e^{D/p}$.  For any $0<c_-<c<c_+$, eventually
\[
 \left(\frac{c_-}{\log n}\right)^{pn}
 \le b_n\le
 \left(\frac{c_+}{\log n}\right)^{pn}.
\]
Write $T=t^{1/p}$.  For a lower bound, fix $0<d<c_-$ and select
$n=\lfloor e^{dT}\rfloor$.  The logarithm of this single term satisfies
\[
 \log(b_nt^n)\ge pn\log(c_-/d)+o(n).
\]
It follows that $\liminf\log\log B(t)/T\ge d$.  Let $d\uparrow c$ through
admissible choices of $c_-$.

For an upper bound choose $d>c_+$ and $N=\lceil e^{dT}\rceil$.
When $n>N$, the upper bound gives
$b_nt^n\le(c_+/d)^{pn}$, up to an immaterial adjustment at the cutoff.
The tail is bounded by a convergent geometric series.  Since the entire
coefficient sequence has bounded $n$th roots, there is $K\ge1$ such that
$b_n\le K^n$ for all $n$.  Therefore the first $N$ terms are bounded by
$N\max(1,Kt)^N$.  Taking two logarithms yields
\[
 \log\log B(t)\le dT+O(\log\log(t+e)).
\]
Let $d\downarrow c$ through choices of $c_+$.  Finite initial terms do not
change either estimate.
\end{proof}

\begin{proof}[Proof of Theorem~\ref{thm:borelintro}]
Take $\beta=0$ in Theorem~\ref{thm:regular} and put $s=p-1$.
After subtracting $s\log n!$, it gives
\[
 \log\frac{u_n}{(n!)^s}
 =-p n\log\log n+
   \left[\log a+p\log\frac p{p-1}\right]n+o(n).
\]
Lemma~\ref{lem:entiregrowth} gives \eqref{eq:doublelog}.
For $p=2$ and $a=1$, the resulting growth is faster than
$\exp(Kt)$ for every fixed $K$.  In fact
$\log B_1(t)=\exp((2+o(1))\sqrt t)$, which dominates $t/q$ for every
$q>0$.  Thus the positive integrand in
\[
 \frac1q\int_0^\infty e^{-t/q}B_1(t)\,\dd t
\]
eventually tends to infinity, so the integral diverges.
\end{proof}

\subsection{The standard Gamma normalization}
For general $s=p-1$, another common normalization is
\[
 \widetilde B_s(t)=\sum_{n\ge1}\frac{u_n}{\Gamma(1+sn)}t^n.
\]
Stirling changes the constant $D$ in
Lemma~\ref{lem:entiregrowth} by $-s\log s$.  Consequently
\begin{equation}
 \lim_{t\to+\infty}
 \frac{\log\log\widetilde B_s(t)}{t^{1/p}}
 =a^{1/p}\frac p s\,s^{-s/p}.
 \label{eq:gammaborelgrowth}
\end{equation}
It follows that $\widetilde B_s$ grows faster on the positive ray than
$\exp(Kt^d)$ for every finite $d,K>0$.  In particular it fails the growth
requirement for the usual order-$1/s$ Borel--Laplace summation procedure.
This is a directional conclusion.  No corresponding assertion about the
negative ray is proved here.

\subsection{Why positivity makes the obstruction especially transparent}
For the quadratic case, Tonelli gives a second proof of the failure:
if the displayed positive-ray Laplace integral were finite, then
\[
 \frac1q\int_0^\infty e^{-t/q}
       \sum_{n\ge1}\frac{u_n t^n}{n!}\,\dd t
 =\sum_{n\ge1}u_nq^n
\]
would be finite.  The convergence criterion says the right-hand side is
infinite.  The explicit double-logarithmic asymptotic is stronger: it
identifies how the growth condition fails rather than merely proving a
contradiction.

\begin{remark}[What ``entire'' does not settle]
The Borel transform here has no finite-plane singularities, but has very
large growth at infinity.  We therefore do not infer ordinary summability
from endless analytic continuation alone, nor do we assign finite Borel
singularities or Stokes constants without an additional argument.  The
separation between analytic continuation and Laplace growth is essential
in the summability framework \cite{sauzin}.
\end{remark}

\section{Sectorial analytic realization without slope bounds}
\label{sec:sector}

The formal series may have zero convergence radius and may fail every
finite Gevrey bound.  Nevertheless its defining equation has a literal
analytic realization in suitable directions.  The reason is a different
estimate: in the left half-plane, $e^{\lambda_j u}$ suppresses rather than
amplifies large positive slopes.

\begin{theorem}[Quantitative left-sector construction]
\label{thm:sector}
Let $\lambda_j\ge0$ be arbitrary and fix $\eta\in(0,1)$.  Set
\[
 K=2e\lambda_1+2.
\]
Choose $r_0>0$ sufficiently small that
\begin{equation}
 r_0\le\frac14,\qquad
 Kr_0\le\min(1,\eta/2),\qquad
 2\lambda_1r_0\le1,
 \label{eq:sectorradius1}
\end{equation}
and
\begin{equation}
 \kappa_0:=r_0\left(\lambda_1+
              \frac{2}{e\eta(1-r_0)}\right)\le\frac12.
 \label{eq:sectorradius2}
\end{equation}
For every $q\in S_{\eta,r_0}$, the map $u\mapsto\PhiL(q,u)$ is a
contraction of the closed disk
\[
 D_q=\{u:|u-q|\le K|q|^2\}
\]
into itself.  Its fixed point $U(q)$ is holomorphic in $q$, satisfies
$U(q)=q+O(q^2)$, and is the unique solution with this latter behavior on
sufficiently small proper subsectors.

Let $U_N(q)$ be the fixed point in the same disk for the first $N$ actions.
Then, for $r=|q|$,
\begin{equation}
 \boxed{\qquad
 |U(q)-U_N(q)|\le
 \frac{r^{N+1}}{(1-r)(1-\kappa_0)}.
 \qquad}
 \label{eq:sectorerror}
\end{equation}
For each fixed $N$, $U_N$ agrees near the origin in the sector with the
analytic germ defined by \eqref{eq:finitekernel}.  The function $U$ has
Poincar\'e expansion $\Uhat$.
\end{theorem}

\begin{proof}
Write $r=|q|$.  If $u\in D_q$, the choices above imply
\[
 |u|\le2r,\qquad \Re u\le-\eta r/2<0.
\]
The infinite sum defining $\PhiL$ is locally normally convergent on
$|q|<1$, $\Re u<0$, since $|q^je^{\lambda_ju}|\le|q|^j$.
It is therefore jointly holomorphic there.

Separate the first action from the rest.  The bound
$|e^z-1|\le|z|e^{|z|}$ gives
\begin{align*}
 |\PhiL(q,u)-q|
 &\le r|e^{\lambda_1u}-1|+\sum_{j\ge2}r^j
 \\
 &\le2e\lambda_1r^2+\frac{r^2}{1-r}
 \le Kr^2.
\end{align*}
Thus the disk is invariant.  For $a>0$ and $t\ge0$,
$t e^{-at}\le1/(ea)$.  Using this with $a=\eta r/2$ gives
\begin{align*}
 |\partial_u\PhiL(q,u)|
 &\le\lambda_1r+
       \sum_{j\ge2}r^j\lambda_j e^{-\lambda_j\eta r/2}
 \\
 &\le\lambda_1r+
       \frac{2}{e\eta r}\frac{r^2}{1-r}
 \le\kappa_0.
\end{align*}
The same estimates hold after truncating the sum at any $N\ge1$.
Banach's contraction theorem supplies the fixed points.  Picard iterates
starting at $q$ are holomorphic on the sector and converge uniformly on
compact subsets, so their limit is holomorphic.

Any other solution $V=q+O(q^2)$ in a proper left subsector has
$|V|\le2r$ and $\Re V\le-\eta r/2$ after reducing $r_0$ if necessary.
The invariance estimate then places $V=\PhiL(q,V)$ in $D_q$, where
uniqueness has already been proved.

Subtract the infinite and finite fixed-point equations.  The segment
joining $U$ and $U_N$ lies in $D_q$, so
\[
 |U-U_N|\le\kappa_0|U-U_N|+
       \sum_{j>N}r^j e^{\lambda_j\Re U}
 \le\kappa_0|U-U_N|+\frac{r^{N+1}}{1-r}.
\]
This proves \eqref{eq:sectorerror}.

For a fixed $N$, the analytic implicit-function theorem at $(0,0)$ gives
the finite-action germ.  It has expansion $q+O(q^2)$ and hence coincides
with the disk fixed point near the origin in the sector.  By
\eqref{eq:finiteagreement}, its first $N$ coefficients are $u_1,\ldots,u_N$.
For that fixed $N$,
\[
 U_N(q)=\sum_{n=1}^Nu_nq^n+O_N(q^{N+1}).
\]
Combining this with \eqref{eq:sectorerror} proves the Poincar\'e assertion.
The constants hidden in $O_N$ may depend on $N$; no uniform Gevrey
remainder estimate has been used or claimed.
\end{proof}

\subsection{The direction of approach is substantive}
If $\sup_j\lambda_j/j=\infty$, no branch satisfying $U(q)/q\to1$ as
$q\downarrow0$ on the positive real axis can solve the literal identity.
Indeed, then $\Re U(q)>0$ for sufficiently small positive $q$, and the
terms $q^je^{\lambda_jU(q)}$ fail to tend to zero along a subsequence.
By contrast, Theorem~\ref{thm:sector} gives a real negative solution when
$q$ is real negative, because the equation and the invariant disk are
preserved by complex conjugation.

Under $q=e^{-x}$, negative $q$ corresponds to a complex horizontal line
with $\Im x$ an odd multiple of $\pi$, rather than to real $x\to+\infty$.
Thus the left-sector construction is a complex realization of the same
formal exponential transseries, not a positive-real realization hidden
by a change of notation.

\subsection{A practical implication for summation algorithms}
The sequence $U_N(q)$, obtained by retaining primitive actions and solving
a finite implicit equation, converges geometrically on these sectors with
the explicit certificate \eqref{eq:sectorerror}.  The Taylor partial sums
$\sum_{n\le N}u_nq^n$ may diverge at every nonzero $q$.  These are different
approximation schemes.  An implementation should record which scheme it
uses, rather than labeling both merely as ``truncation of the transseries.''

\section{Transfer to genuine compositional inversion}
\label{sec:inverse}

Because $u_1=1$, the formal solution has a unique inverse
$\Qhat\in q+q^2\R[[q]]$ under composition.  In particular,
\begin{equation}
 \Qhat(\Uhat(q))=q,
 \qquad
 u=\sum_{j\ge1}\Qhat(u)^j e^{\lambda_j u}
 \label{eq:inverseidentity}
\end{equation}
as formal identities.  The following classical stability fact is proved
here to keep the application independent of analytic realization choices.

\begin{proposition}[Gevrey stability under formal reversion]
\label{prop:reversion}
If $\widehat f(q)=q+\sum_{n\ge2}f_nq^n$ and $s\ge0$, then
\[
 \widehat f\in\Gev_s
 \quad\Longleftrightarrow\quad
 \widehat f^{\langle-1\rangle}\in\Gev_s.
\]
\end{proposition}

\begin{proof}
For $s=0$ this is the ordinary analytic inverse-function theorem in both
directions.  We give a coefficient proof valid also for $s>0$.
Assume $|f_n|\le CA^n(n!)^s$, increasing $A,C$ to at least one.
Write $\widehat f=q(1+h(q))$, where
$h(q)=\sum_{m\ge1}f_{m+1}q^m$.  Lagrange inversion gives, for $n\ge2$,
\begin{align*}
 [q^n]\widehat f^{\langle-1\rangle}
 &=\frac1n[z^{n-1}](1+h(z))^{-n}\\
 &=\frac1n\sum_{k=1}^{n-1}(-1)^k\binom{n+k-1}{k}
       \sum_{\substack{m_1+\cdots+m_k=n-1\\m_i\ge1}}
            \prod_i f_{m_i+1}.
\end{align*}
For positive integers $m_i$ with sum $n-1$,
\begin{equation}
 \prod_{i=1}^k(m_i+1)!
 \le2^{k-1}(n-k+1)!\le2^n n!.
 \label{eq:factorialmerge}
\end{equation}
One way to see the first inequality is to transfer a unit from a smaller
factorial argument greater than two to a larger one; this cannot decrease
the product, and it leaves all but one argument equal to two.

The binomial coefficient in the reversion formula is at most $4^n$,
the number of compositions is at most $2^n$, and
$\prod_i CA^{m_i+1}\le C^nA^{2n}$.  Equation~\eqref{eq:factorialmerge}
therefore bounds the inverse coefficient by
$B^n(n!)^s$ for a fixed $B>0$, after absorbing all exponential factors.
This proves one implication.  Apply it again to the inverse series to
obtain the other.
\end{proof}

\begin{corollary}[The classification survives inversion]
\label{cor:inverse}
The inverse $\Qhat$ has positive convergence radius exactly when
$\lambda_j=O(j)$.  For each $s>0$,
\[
 \Qhat\in\Gev_s
 \quad\Longleftrightarrow\quad
 \lambda_j=O((j\log(j+1))^{s+1}).
\]
Its least Gevrey index is the value in \eqref{eq:index}.
\end{corollary}

\begin{remark}
This corollary transfers \emph{class membership}, not the numerical type
$T_s$.  The inverse coefficients need not be nonnegative.  We have not
proved that $T_s(\Qhat)=T_s(\Uhat)$, nor that the concentration theorem
transfers through the cancellations in reversion.
\end{remark}

\section{Examples, benchmarks, and sharp boundaries}
\label{sec:examples}

\subsection{Analytic benchmark kernels}

When all slopes vanish, $U=q/(1-q)$ has radius one.  If $\lambda_j=1$
for every $j$, then
\[
 U=\frac{q}{1-q}e^U,
 \qquad U=-W_0\!\left(-\frac q{1-q}\right).
\]
Its positive critical point is $U=1$, $q=1/(1+e)$; the positive
supersolution criterion and the nonzero quadratic derivative at the
critical point show that this is its radius of convergence.

For $\lambda_j=j$, the kernel is geometric in $qe^U$, and
\begin{equation}
 U=\frac{qe^U}{1-qe^U}
 \quad\Longleftrightarrow\quad U=q(1+U)e^U.
 \label{eq:linearbenchmark}
\end{equation}
Writing $q=Ue^{-U}/(1+U)$, the positive maximum occurs at
\[
 \tau=\frac{\sqrt5-1}{2},\qquad
 R=\frac{\tau e^{-\tau}}{1+\tau}=\tau^2e^{-\tau}.
\]
The branch at the origin reaches this simple critical point.  Local
analytic inversion gives
\begin{equation}
 U(q)=\tau-
 \sqrt{\frac{2(1+\tau)}{3+\tau}}\sqrt{1-q/R}
 +O(1-q/R).
 \label{eq:linearsqrt}
\end{equation}
This is a familiar analytic implicit-function singularity; the new
phenomena begin when the countable slopes outgrow their actions.

\subsection{Quadratic feedback}
For $\lambda_j=j^2$,
\begin{align*}
 \Uhat(q)={}&q+2q^2+\frac{15}{2}q^3+\frac{107}{3}q^4
 +\frac{1555}{8}q^5+\frac{17362}{15}q^6\\
 &+\frac{5288311}{720}q^7+\frac{15414767}{315}q^8+\cdots.
\end{align*}
It has radius zero, is Gevrey-$1$ of type zero, and is not Gevrey-$s$
for any $s<1$.  Its coefficients satisfy
\begin{equation}
 \log u_n=n\log n-2n\log\log n+(2\log2-1)n+o(n).
 \label{eq:quadraticasymptotic}
\end{equation}
Thus ``Gevrey-$1$ of type zero'' is compatible with divergence: the
normalization by $n!$ is stronger than the normalization by an exponential.
The ordinary Borel transform is entire but has the positive-ray growth
proved in Section~\ref{sec:borel}.

\subsection{The critical logarithmic correction}
For $\lambda_j=a(j\log(j+1))^2$, the exact type identity gives
\[
 T_1(\Uhat)=4a,
 \qquad R_{\mathrm{Borel}}=\frac1{4a}.
\]
This finite Borel radius is obtained without a full multiplicative
asymptotic for the coefficients.  Increasing the logarithmic exponent to
$2+\varepsilon$ destroys Gevrey-$1$ membership; decreasing it below two
makes the ordinary Borel transform entire.

\begin{center}
\small
\begin{tabular}{@{}p{0.30\textwidth}p{0.19\textwidth}p{0.39\textwidth}@{}}
\toprule
Slope sequence & Convergence & Gevrey conclusion\\
\midrule
$\lambda_j=O(j)$ & Positive radius & $\Gev_0$\\
$j\log(j+1)$ & Radius zero & Every $\Gev_s$, $s>0$; not $\Gev_0$\\
$j^2$ & Radius zero & Critical class $\Gev_1$, type zero\\
$a(j\log(j+1))^2$ & Radius zero & Critical class $\Gev_1$, type $4a$\\
$j^2(\log(j+1))^3$ & Radius zero & Every $\Gev_s$, $s>1$; not $\Gev_1$\\
$j^3$ & Radius zero & Critical class $\Gev_2$, type zero\\
$e^{\sqrt j}$ & Radius zero & No finite Gevrey class\\
\bottomrule
\end{tabular}
\end{center}
All rows nevertheless admit the left-sector realization of
Theorem~\ref{thm:sector}.  The last row is a particularly strong warning
against inferring coefficient regularity from sectorial existence alone.

\subsection{Why a complete black-box decision procedure is impossible}
The sharp criteria are mathematical equivalences, not a claim that their
asymptotic conditions can be decided from an arbitrary program for
$\lambda_j$.

\begin{proposition}
\label{prop:undecidable}
There is no algorithm that takes an arbitrary program computing a
nonnegative integer slope sequence and always decides whether its formal
solution is convergent.  For any fixed $s>0$, there is likewise no such
algorithm deciding membership in $\Gev_s$.
\end{proposition}

\begin{proof}
Given a machine $M$, define a total computable sequence by running $M$
for $j$ steps and setting
\[
 \lambda_j=\begin{cases}
 j^2,&M\text{ has not halted within those steps},\\
 0,&M\text{ has halted within those steps}.
 \end{cases}
\]
If $M$ halts, only finitely many slopes are nonzero and the solution is
convergent.  If $M$ never halts, the sequence is $j^2$ and the solution is
divergent.  A proposed decision procedure would decide the halting problem.
For a fixed positive $s$, replace $j^2$ by $j^m$ for an integer $m>s+1$.
The nonhalting case is not in $\Gev_s$, while the halting case remains
convergent.  The same reduction applies.
\end{proof}

The appropriate implementation goal is consequently a sound certificate
system for structured input families, not an unconditional classifier for
arbitrary programs.  This reduction is a standard use of undecidability;
it is not presented as a new theorem of computability theory.

\section{Reproducible coefficient checks and numerical diagnostics}
\label{sec:computations}

The supplementary script \texttt{code/verify.py} uses exact integers for
coefficient generation when the slopes are nonnegative integers.  It also
uses an independent partition enumeration to check the Lagrange formula.
No numerical experiment is used as a substitute for a proof above.

\subsection{An exact recurrence}
Put
\[
 A_n=n!u_n,
 \qquad E_{j,k}=k![q^k]e^{\lambda_j\Uhat(q)},\qquad E_{j,0}=1.
\]
Differentiating the exponential formally gives
\begin{equation}
 E_{j,k}=\lambda_j\sum_{r=1}^k
           \binom{k-1}{r-1}A_rE_{j,k-r}.
 \label{eq:exactrecurrence1}
\end{equation}
The original fixed-point identity gives
\begin{equation}
 A_n=\sum_{j=1}^n\frac{n!}{(n-j)!}E_{j,n-j}.
 \label{eq:exactrecurrence2}
\end{equation}
All quantities needed on the right have smaller index than $n$, except
for the already known constant $E_{n,0}=1$.  Integer slopes therefore
produce integer $A_n$, even when $u_n$ is not integral.  The supplied
implementation takes $O(N^3)$ integer arithmetic operations and
$O(N^2)$ stored integers through order $N$; this is an arithmetic-operation
count, not a bit-complexity bound.

The executed checks were as follows.  For $\lambda_j=j^2$, coefficients
through order $120$ were computed exactly.  The first twelve were checked
against direct partition enumeration.  Twelve zero-slope coefficients
were checked against $q/(1-q)$, and twelve linear-slope coefficients
against an independent Lagrange formula for \eqref{eq:linearbenchmark}.
All thirty-six comparisons passed.  The exact values and the verification
record are included in \texttt{data/}.

\subsection{A rigorous sandwich, evaluated approximately}
For the quadratic model and $k=n-r+1$, convexity gives the exact maximum
\[
 \sum_{i=1}^k j_i^2\le r^2+k-1.
\]
Consequently
\begin{align}
 u_n&\ge
 \max_{2\le r\le n}
 \frac{(r^2+n-r)^{n-r}}{(n-r)!},
 \label{eq:quadraticlower}\\
 u_n&\le\sum_{r=1}^n\binom{n-1}{r-1}
       \frac{(r^2+n-r)^{n-r}}{(n-r+1)!}.
 \label{eq:quadraticupper}
\end{align}
For numerical work the all-unit contribution can also be included in the
lower maximum, with its correct denominator $n!$.  The mathematical
inequalities are rigorous.  The decimal evaluations below use ordinary
floating-point logarithms, not directed rounding, and are diagnostics
rather than interval certificates.

Define the normalized logarithm
\[
 H_n=\frac{\log u_n}{n}-\log n+2\log\log n.
\]
Theorem~\ref{thm:regular} proves $H_n\to2\log2-1\approx0.38629436$.
Applying the normalization to the lower and upper bounds gives:
\begin{center}
\begin{tabular}{@{}rrrrr@{}}
\toprule
$n$ & Lower value & Upper value & Lower-bound maximizing $r$ & $r\log n/n$\\
\midrule
$100$ & $1.017199$ & $1.626109$ & $38$ & $1.749965$\\
$1\,000$ & $1.059262$ & $1.666026$ & $292$ & $2.017065$\\
$10\,000$ & $1.031046$ & $1.587396$ & $2\,332$ & $2.147851$\\
$100\,000$ & $0.990011$ & $1.493908$ & $19\,183$ & $2.208524$\\
$1\,000\,000$ & $0.949010$ & $1.406654$ & $161\,926$ & $2.237090$\\
\bottomrule
\end{tabular}
\end{center}
These values make an important caution visible: convergence of the
normalized logarithm and of the leading saddle scale is slow.  A moderate
coefficient table would not reliably identify the limiting constant or
justify a full coefficient asymptotic.  The table is not evidence that the
finite-$n$ maximizing action is already close to its limiting value two.

\subsection{Reproduction commands}
From the archive root, run
\begin{quote}\ttfamily\small
python -m pip install -r requirements.txt\\
python code/verify.py --order 120\\
pdflatex -interaction=nonstopmode -halt-on-error article.tex\\
pdflatex -interaction=nonstopmode -halt-on-error article.tex
\end{quote}
The exact computation uses only the Python standard library.  NumPy and
SciPy are used for the large-$n$ floating-point envelope tables.
The article has no external figure or bibliography-file dependency.

\section{A focused further-research program}
\label{sec:future}

The following questions extend proved results rather than presupposing
unproved summability or novelty claims.  Their priority refers to their
usefulness for the repository's transseries-inversion program.

\subsection{Directional summability of the quadratic model}
Determine the growth of $B_1(t)$ on nonpositive rays and in sectors, not
only on the positive axis.  Is the left-sector solution of
Theorem~\ref{thm:sector} the ordinary Borel sum of $\Uhat$ in direction
$\pi$?  A successful proof must establish the required sectorial Borel
growth and compare the Laplace sum with the literal fixed point.
The coefficient Gevrey bound and Poincar\'e realization alone do not imply
this comparison.  Cancellation in $B_1(-t)$ is the central new issue.
% ed. (2026-09-29): answered by later sibling packages (negative-ray, natural-boundaries, sectorial-summability).
\begin{ednote}
Later packages in \nolinkurl{Analysis/Transseries/docs/series-and-transseries/}
answer this question in three senses; none of them has been independently
reviewed. (i)~\emph{Fine sense, answered positively.}
\nolinkurl{Negative_Ray_Summation_Exponential_Feedback/} proves
(its Theorem~\texttt{thm:threshold}) that, for these unit amplitudes, $\Uhat$ and
$\Qhat$ are finely (fixed-width half-strip) Borel--Laplace summable in
direction $\pi$ exactly when $\lambda_j=O((j\log j)^2)$, which includes the
quadratic model, with sums that satisfy the literal kernel near the
negative axis; its Theorem~\texttt{thm:borel} continues the inverse Borel transform
to $|\zeta|+\Re\zeta<1/(2A)$, $A=\limsup\lambda_j/(j\log j)^2$.
(ii)~\emph{Angular sense, answered negatively for $\lambda_j=j^2$.}
\nolinkurl{Natural_Boundaries_Quadratic_Exponential_Feedback/} proves
(its Theorem~\texttt{thm:main}) that the Borel transform of $\Qhat$ is entire and
exponentially bounded on every tube around the negative ray but in no
wedge around it, so neither $\Qhat$ nor $\Uhat$ is $1$-summable in an open
arc of directions containing $\pi$. (iii)~\emph{Orders $k>1$.}
\nolinkurl{Sharp_Negative_Direction_Summability_Exponential_Feedback/}
(its Theorems~\texttt{thm:headline} and \texttt{thm:ksum}) proves $k$-summability in
direction $\pi$, with sum the canonical left-sectorial solution, exactly
when $\lambda_j=O((j\log(j+1))^{1+1/k})$; it leaves $k=1$ open.
\end{ednote}

\subsection{Acceleration and the choice of analytic equation}
The law $\log\log B_1(t)\sim2\sqrt t$ suggests that ordinary Laplace
summation is not the appropriate positive-direction operation.  Develop an
accelerated transform adapted to this growth and state explicitly which
analytic continuation of the countable kernel it respects.  A literal
positive-real solution of \eqref{eq:kernel} is impossible when
$\lambda_j/j$ is unbounded, so an accelerated value cannot be claimed to
solve that divergent literal sum without specifying an additional
continuation or regularization rule.  Existence, uniqueness, and dependence
on that choice are separate research questions.

\subsection{Beyond logarithmic coefficient asymptotics}
Refine \eqref{eq:logasymptotic} to include its sublinear terms and eventually
a full multiplicative equivalent.  The exact composition entropy, the
small nonunit actions, and the shift of the large-action saddle all become
relevant below order $n$.  Establish a local limit theorem for $K$ and
$J_{\max}$ with a correctly centered fluctuation law.  The exponential
concentration proved here supplies localization, but not the necessary
subleading analysis.
% ed. (2026-09-29): answered for power slopes by the finite-core and microscopic-condensation packages.
\begin{ednote}
Answered for two slope classes by later packages. For $\lambda_j=aj^p$,
$p>1$, after finitely many nonnegative changes,
\nolinkurl{Finite_Core_Universality_Exponential_Feedback/} proves a
multiplicative saddle equivalent with its determinant prefactor
(Theorem~\texttt{thm:saddle} there), a conditional local Gaussian law and an
unconditional central limit theorem for the largest action
(Theorem~\texttt{thm:clt}), and an explicit first correction for $p>3/2$
(Proposition~\texttt{prop:first-correction}). For
$\lambda_j=j^2+\alpha j+\eta j/\log j+o(j/\log j)$,
\nolinkurl{Microscopic_Condensation_Exponential_Feedback/} proves a full
multiplicative equivalent (Theorem~\texttt{thm:coefficient}), which for eventually
exact $j^2$ is the $p=2$, $a=1$ case of the finite-core first-correction
formula, joint total-variation Gaussian--Poisson laws for the largest
action and the number of twos, a two-dimensional local limit, and a
central limit theorem for $K$ (Theorem~\texttt{thm:probability},
Corollaries~\texttt{cor:local} and \texttt{cor:K}).
\nolinkurl{Poisson_Layers_Finite_Core_Boundary/} adds the first omitted
action as an explicit Poisson layer. Slopes $aj^p(\log j)^\beta$ with
$\beta\ne0$ are not treated in these packages.
\end{ednote}

\subsection{The near-linear transition}
Analyze slopes such as $j(\log j)^\beta$ with $\beta>0$ and
$jL(j)$ with a slowly varying unbounded $L$.  The present classification
already proves divergence together with membership in every positive
Gevrey class.  What are the sharp coefficient asymptotics and the natural
summation scales in this zero-index, nonconvergent regime?  The saddle
balance changes because the $(p-1)n\log n$ term disappears at $p=1$.
% ed. (2026-09-29): answered at n-th-root precision by the near-linear package; negative direction by two others.
\begin{ednote}
Answered at $n$-th-root precision by
\nolinkurl{Near_Linear_Boundary_Exponential_Feedback/} for
$\lambda_j=jL(j)$ with $L$ positive, nondecreasing, $C^1$ and unbounded,
of elasticity tending to zero, with $xL(x)$ convex, and with exponentially
controlled positive amplitudes: $u_n^{1/n}\sim c_1e^{b_n}$ where
$b_ne^{b_n}=L(n/(1+b_n))$ (its Theorem~\texttt{thm:main}), with an explicit law for
$j(\log j)^\beta$; every factorial-normalized Borel transform is entire but
grows too fast for Laplace summation in the positive direction
(Theorem~\texttt{thm:borel}, Corollary~\texttt{cor:no-laplace}). A multiplicative equivalent
for $u_n$ is not claimed there. In the negative direction this class has
$A=0$ in the negative-ray package (fine summability and an entire inverse
Borel transform, for unit amplitudes), and the sectorial-summability
package gives strong Gevrey-$s$ remainders for every $s>0$ and
$k$-summability for every $k>1$.
\end{ednote}

\subsection{Exact type under compositional inversion}
For $\Qhat=\Uhat^{\langle-1\rangle}$, determine whether the exact type
$T_s$ is preserved and under which hypotheses.  We have proved only class
membership and equality of the least Gevrey index.  A proof of a sharper
claim must control cancellations in the reversion formula rather than
import the positive-weight argument unchanged.
% ed. (2026-09-29): answered by the weighted-type package; multiplicative inverse laws by the finite-core and batch-49 packages.
\begin{ednote}
Answered by \nolinkurl{Sharp_Weighted_Type_Formal_Reversion/}. Its
kernel-independent zero-loss theorem (Theorem~\texttt{thm:zeroloss} there) shows
that tangent-to-identity reversion preserves the weighted type for
log-convex weights with $M_n^{1/n}\to\infty$; combined with
\eqref{eq:exacttype}, which it imports from this article and restates but
does not re-prove, it gives $T_s(\Qhat)=T_s(\Uhat)$ for every $s>0$
(Theorem~\texttt{thm:inversefeedback}), the inverse Borel radius $1/(4A)$
(Corollary~\texttt{cor:borelradius}) and a refined inverse type for
$\lambda_j\sim aj^p(\log j)^\beta$ (Theorem~\texttt{thm:refinedfeedback}).
Multiplicative inverse asymptotics are proved for $\lambda_j=aj^p$,
$p>2$, in the finite-core package (its Theorem~\texttt{thm:explicit}); at $p=2$
that package states a conjecture (\texttt{conj:quadratic-inverse}), of which
\nolinkurl{Quadratic_Exponential_Feedback_After_Reversion/} and
\nolinkurl{Signed_Quadratic_Feedback_Inversion/} claim independent,
not yet reviewed, proofs.
\end{ednote}

\subsection{Amplitudes, general actions, and multiple scales}
Replace the unit amplitudes and integer actions by
\[
 U=\sum_j c_j q^{a_j}e^{\lambda_j U},
 \qquad c_j\ge0,
\]
first with positive integer $a_j$ and then with well-based real actions.
Determine the analogues of the exact type identity in terms of the
interaction between $c_j$, $a_j$, and $\lambda_j$.  The proof here uses the
availability of arbitrarily many action-one factors in an essential way.
For true multilevel transseries, one must also state which filtration is
used for formal summability and which parameter controls the analytic
normalization.  A weight estimate in the wrong filtration is not an
analytic certificate.
% ed. (2026-09-29): partly answered for special amplitude classes (near-linear, critical Hahn, natural-boundaries packages).
\begin{ednote}
Partly answered for special amplitude classes. The near-linear package
allows positive amplitudes with $e^{-D(j-1)}\le c_j/c_1^j\le e^{D(j-1)}$.
\nolinkurl{Critical_Hahn_Transseries_Beyond_Finite_Action_Folds/}
treats $U=\beta\sum_ja_jq^je^{jU}$ with $a_j=aj^{-1-\alpha}$ outside a
finite set, $1<\alpha<2$, at the critical coupling (a convergent Hahn
chart, stable-density sector asymptotics and action budgets); it does not
cite this question. The natural-boundaries package relates rational
amplitude generating functions with eventually quadratic slopes to angular
summability (its Theorem~\texttt{thm:rational-main}). None of them states an
analogue of \eqref{eq:exacttype} in terms of $c_j$, $a_j$ and $\lambda_j$.
\end{ednote}

\subsection{Signed and vector-valued feedback}
For complex amplitudes or slopes, positivity no longer identifies the
absolute majorant with the actual coefficient growth.  Seek phase or
noncancellation conditions under which a sharp lower bound survives.
For vector systems, investigate whether a scalar slope budget should be
replaced by spectral radii of nonnegative feedback matrices.  An adequate
theorem should include both a finite-dimensional reduction and a
countable-tail hypothesis, rather than silently use a finite-dimensional
implicit-function theorem on an infinite system.

\subsection{Uniform Gevrey remainder bounds for the sectorial solution}
Our left-sector construction has an explicit primitive-action truncation
error, while its Poincar\'e remainder constants are allowed to depend
arbitrarily on order.  Determine when the actual solution satisfies
uniform estimates of the form
\[
 \left|U(q)-\sum_{n<N}u_nq^n\right|
 \le CA^N(N!)^s|q|^N
\]
on proper sectors.  The sharp coefficient criterion is necessary for
such a statement, but its sufficiency for this particular realization
requires a new argument.  This is the most direct bridge to the Gevrey
implicit-function and summability literature.
% ed. (2026-09-29): answered by the sectorial-summability package; quadratic inverse remainder in the negative-ray package.
\begin{ednote}
Answered by
\nolinkurl{Sharp_Negative_Direction_Summability_Exponential_Feedback/}
(its Theorem~\texttt{thm:headline}): for every $s>0$ the displayed estimate holds on
every proper left sector, for the canonical left-sectorial solution that
package reconstructs (Proposition~\texttt{prop:existence}), exactly when
$\lambda_j=O((j\log(j+1))^{s+1})$, and likewise for the compositional
inverse (Theorem~\texttt{thm:inverse}). Its constants depend on the sector; no
estimate up to the boundary rays is claimed. For $\lambda_j=j^2$ the
negative-ray package proves, for $P=e^uQ$, a uniform remainder with an
explicit majorant on a closed left half-disc (its Theorem~\texttt{thm:remainder}).
\end{ednote}

\subsection{Formalization and certified computation}
A useful Lean development can be staged without first formalizing all of
transseries theory.  The finite coefficient formula, coefficientwise
monotonicity, the one-large-action lower bound, and the convex composition
bound are natural initial targets.  A second stage should formalize the
logarithmic envelope and its limit estimates.  A third stage can address
Tonelli, normal convergence, and sectorial contraction.

For computation, replace the illustrative logarithmic tables by directed
interval arithmetic, and attach proof objects to verified slope bounds.
The positive and left-sector certificates should remain separate in the
API.  A result should record whether it establishes formal support,
coefficient Gevrey growth, an analytic germ, a sectorial solution,
a uniform asymptotic expansion, or Borel summability.  These are genuinely
different outputs.

\section{Conclusions}
\label{sec:conclusion}

The countable kernel \eqref{eq:kernel} is simple enough to admit exact
coefficient formulas yet rich enough to distinguish several regularity
notions that can otherwise be conflated in transseries inversion.
Its formal solution always exists.  Linear slope growth is exactly the
threshold for ordinary convergence.  For every positive Gevrey order,
the sharp threshold is instead $(j\log j)^{s+1}$, and the type is given
by the explicit constant in \eqref{eq:exacttype}.  The threshold survives
formal compositional inversion at the level of class membership.

For regularly varying superlinear slopes, the coefficients are generated
at leading exponential precision by one large action together with many
unit-action factors.  This gives both the term of order $n$ in
$\log u_n$ and an exponential concentration theorem.  In the quadratic
case, the Borel transform is entire but too large for positive-direction
Laplace summation.  Finally, suppression of large slopes in the left
half-plane gives a constructive analytic realization for every
nonnegative slope sequence, with a uniform action-truncation certificate.

The resulting improvement to an inversion workflow is concrete: formal
well-basedness should be accompanied by an explicit slope budget before
analytic or Gevrey regularity is asserted.  When a positive majorant fails,
a direction-sensitive analytic construction may still be available, but
its output must not be confused with a convergent Taylor series or a
proved Borel sum.

\appendix
\section{Dependency map and proof-audit checklist}
\label{app:audit}

\begin{center}
\small
\begin{tabular}{@{}p{0.34\textwidth}p{0.53\textwidth}@{}}
\toprule
Result & Main dependencies and boundary\\
\midrule
Formal fixed point & Finite coefficient support and $q$-adic contraction; no analytic assumption.\\
Positive coefficient formula & Classical formal Lagrange inversion; $k\le n$ makes specialization finite.\\
Convergence equivalence & Positivity, Tonelli, and a product-neighborhood analytic argument.\\
Logarithmic asymptotic & Eventual convexity, one-large-action lower bound, sharp envelope lemma.\\
Exact Gevrey type & Logarithmic asymptotic for an upper comparison sequence; sparse-subsequence lower bound.\\
Concentration & Strict concavity of the envelope and a capped convex composition bound.\\
Entire Borel growth & Positive coefficients and an elementary maximum-term estimate.\\
Sectorial realization & Suppression for $\Re u<0$, disk contraction, normal convergence.\\
Poincar\'e expansion & Finite-action germs plus the uniform action-tail bound; no uniform Gevrey remainder claim.\\
Reversion classification & Classical coefficient inversion and factorial bounds; no exact-type transfer claimed.\\
\bottomrule
\end{tabular}
\end{center}

For independent review, the most important points to check are the
localization steps in Lemma~\ref{lem:envelope}, the retention of the
$j/n$ correction in \eqref{eq:typelowercalc}, and the distinction between
primitive-action truncation and Taylor truncation.  The constant in
\eqref{eq:exacttype} is sensitive to the second of these points.
The positivity assumption is used in the necessity direction of the
convergence criterion and in all sharp lower bounds.

\section{Repository integration note}

A natural location for this note is a new research subdirectory adjacent
to the consolidated transseries volume, with its own proof-status record.
No existing repository source or branch was modified in preparing this
archive.  The package contains the article source and compiled PDF,
a verification script, machine-readable coefficient and envelope tables,
and a README explaining reproduction and limitations.

The bibliography distinguishes the repository motivation from external
classical tools.  The sharp classification theorems are derived within
this article; they are not quoted from the repository.  The literature
check located established infinite-colour inversion and Gevrey implicit
function theory, but was not an exhaustive priority search.  Independent
mathematical review and further bibliographic comparison are required
before treating the proposed results as a publication-level novelty claim.

\clearpage
\phantomsection
\addcontentsline{toc}{section}{References}
\begin{thebibliography}{9}

\bibitem{repo}
V. Reshetnikov,
\emph{ProveIt: Transseries---the polynomial--logarithmic calculus and its
inversions}, repository source and README,
commit \texttt{04e06e032dff1966513bfba316966d52db3646b4},
consulted 29 September 2026.
\newblock
\href{https://github.com/VladimirReshetnikov/ProveIt/tree/04e06e032dff1966513bfba316966d52db3646b4/Analysis/FabiusFunction/docs/semi-formalized-research-frontiers/drafts/series-and-transseries/Transseries_And_Inversion}{Pinned transseries directory on GitHub}.
% ed. (2026-09-29): current location of the pinned directory added.
\newblock Editorial note (ProveIt, 2026-09-29): the pinned link above is
valid at that commit; the directory is now
\nolinkurl{Analysis/Transseries/docs/series-and-transseries/Transseries_And_Inversion/}.

\bibitem{vdh}
J. van der Hoeven,
\emph{Transseries and Real Differential Algebra},
Lecture Notes in Mathematics 1888, Springer, 2006.
\newblock Author-hosted text:
\url{https://www.texmacs.org/joris/ln/ln.html}.

\bibitem{gessel}
I. M. Gessel,
Lagrange inversion,
\emph{Journal of Combinatorial Theory, Series A} \textbf{144} (2016),
212--249.
\newblock DOI: \href{https://doi.org/10.1016/j.jcta.2016.06.018}{10.1016/j.jcta.2016.06.018}.
\newblock Author preprint: \url{https://arxiv.org/abs/1609.05988}.

\bibitem{jktcolors}
S. Jansen, T. Kuna, and D. Tsagkarogiannis,
Lagrange inversion and combinatorial species with uncountable color palette,
\emph{Annales Henri Poincar\'e} \textbf{22} (2021), 1499--1534.
\newblock DOI: \href{https://doi.org/10.1007/s00023-020-01013-0}{10.1007/s00023-020-01013-0}.
\newblock Preprint: \url{https://arxiv.org/abs/2008.10862}.

\bibitem{jktvirial}
S. Jansen, T. Kuna, and D. Tsagkarogiannis,
Virial inversion and density functionals,
\emph{Journal of Functional Analysis} \textbf{284}, no. 1 (2023), 109731.
\newblock DOI: \href{https://doi.org/10.1016/j.jfa.2022.109731}{10.1016/j.jfa.2022.109731}.
\newblock Preprint: \url{https://arxiv.org/abs/1906.02322}.

\bibitem{nikolaev}
N. Nikolaev,
Gevrey asymptotic implicit function theorem,
\emph{L'Enseignement Math\'ematique} \textbf{70}, no. 1/2 (2024), 251--282.
\newblock DOI: \href{https://doi.org/10.4171/LEM/1061}{10.4171/LEM/1061}.
\newblock Preprint: \url{https://arxiv.org/abs/2112.08792}.

\bibitem{sauzin}
D. Sauzin,
\emph{Introduction to 1-summability and resurgence},
2014 lecture notes.
\newblock \url{https://arxiv.org/abs/1405.0356}.

\end{thebibliography}
\end{document}
